\chapter{Marco metodológico}

\section{Metodología}
El equipo (tres personas) trabajó con \textbf{Scrum adaptado}: un backlog de producto priorizado con MoSCoW, releases con un resultado de producto
explícito, un incremento demostrable por sprint, criterios de aceptación en Gherkin y un registro escrito de las decisiones de arquitectura
(ADR). Tres prácticas gobernaron el trabajo en la práctica:

\begin{itemize}
  \item \textbf{Lo que ya funciona no se rompe.} Cada cambio mantiene en verde todas las pruebas anteriores; se prefiere extender con
  parámetros opcionales antes que reescribir un comportamiento probado.
  \item \textbf{Un commit por tema, probado en su estado exacto.} El historial de Git (175 commits entre el 9 de septiembre y el 20 de
  noviembre de 2026) es verificable y cuenta la secuencia de decisiones.
  \item \textbf{Evidencia antes que afirmación.} Cada atributo se midió con una herramienta (JUnit, k6, Prometheus, Sonar) y el resultado
  quedó versionado junto al código.
\end{itemize}

\textbf{Herramientas.} Repositorio y control de versiones en GitHub; integración y despliegue continuos con GitHub Actions
(\emph{reusable workflows}); calidad con SonarQube/SonarCloud; modelo de arquitectura con ArchiMate (Archi); documentos en Markdown y
LaTeX; backlog, release plan y ADR versionados en \cod{docs/}.

\textbf{Roles observables en el repositorio} (según los commits): Hever Barrera, núcleo de detección, seguimiento y alertas, gateway,
panel, persistencia y los atributos de calidad del Release~2; María Juliana Rodríguez, integración del inicio de sesión con Google,
pipeline de CI/CD, calidad estática y despliegue en AWS; Kevyn Forero, estructura inicial, integración de back y front con Docker e
instrumentación de observabilidad (Micrometer, Prometheus y Grafana). La sección~\ref{sec:contribuciones} detalla las contribuciones.

\section{Criterios DoR y DoD}
\begin{table}[H]
\centering\small
\caption{Definition of Ready y Definition of Done del proyecto (docs/backlog.md, sección 3).}\label{tab:dordod}
\begin{tabularx}{\textwidth}{YY}
\toprule
\textbf{DoR: una HU entra a desarrollo cuando\ldots} & \textbf{DoD: una HU está terminada cuando\ldots} \\
\midrule
Tiene el formato «Como \ldots quiero \ldots para \ldots» (no es una tarea técnica) &
El código está integrado y compila \\
Tiene al menos un criterio de aceptación en Gherkin con datos concretos &
Sus pruebas automáticas pasan y todas las anteriores siguen pasando \\
Sus dependencias están identificadas (otra HU, cuenta externa, permiso) &
Si toca la interfaz, se verificó en un navegador real \\
Su tamaño es de 8 puntos o menos (si no, se divide) &
Las claves y contraseñas están solo en variables de entorno, nunca en el código, los registros ni el repositorio \\
Si toca datos de personas, se revisó la privacidad (qué sale a terceros) &
Hay evidencia (prueba, medición o demostración) y la documentación (ADR o README) está al día; se entrega en un commit por tema \\
\bottomrule
\end{tabularx}
\end{table}

\section{Planificación y estimación}
\textbf{Estimación.} Las historias se estiman en puntos con la escala 1, 2, 3, 5 y 8 (una historia de más de 8 se divide). El backlog
tiene \textbf{38 historias y 145 puntos}; al cierre del Release~2 hay \textbf{33 hechas (128 puntos, el 88\,\%)}, una parcial (HU-30,
Telegram con lista fija de chats) y cuatro pendientes (HU-20, 28, 31 y 32; 12 puntos). De las 19 historias \emph{Must}, 18 están hechas;
la pendiente es probar con dos cámaras físicas reales (HU-28), que depende de una segunda cámara.

\textbf{Releases por resultado de producto.} El plan (\cod{docs/release-plan.md}) agrupa los sprints por la capacidad que dejan
funcionando, no por tareas técnicas:

\begin{table}[H]
\centering\small
\caption{Plan de releases del producto (ocho sprints lógicos de la capacidad de producto; la Tabla~\ref{tab:sprints} los asigna a los
sprints del curso).}\label{tab:releases}
\begin{tabularx}{\textwidth}{lYY}
\toprule
\textbf{Release} & \textbf{Resultado de producto} & \textbf{Sprints (incremento)} \\
\midrule
R1 «Ver» & El sistema distingue personas y maletas y sabe cuáles se dejaron en reposo & S1 detección; S2 registro de objetos en reposo con foto \\
R2 «Alertar» & Que el personal de seguridad se entere en menos de 5~s de que alguien se llevó un objeto, con evidencia & S3 retiro y dueño; S4 alerta con evidencia; S5 panel en tiempo real y Docker \\
R3 «Operar y avisar» & Usar las cámaras existentes (hasta 2), reducir falsas alarmas y avisar fuera del panel & S6 dos cámaras; S7 reporte y «dueño que volvió»; S8 avisos y acceso con Google \\
\bottomrule
\end{tabularx}
\end{table}

\textbf{Riesgos gestionados.} El plan mantiene una tabla de riesgos con su mitigación y su estado. Los más relevantes: la falsa alarma cuando
el dueño sale de cuadro y vuelve (mitigada: baja a «posible dueño que volvió», media, sin silenciar, con riesgo residual declarado); el
rendimiento con dos cámaras (mitigado: un modelo más liviano con los hilos repartidos); y la privacidad (controlada: no se guarda video y los
avisos externos son solo texto).
